% =====================================================================
%  26.4  Recognition -- reading a statement and spotting the VT signature.
%  Exports: sec:c30-recognition  def:c30-signature
%  WORKED reading only -- NO solvedproblem (in-section solved stays 2: 2127E,
%  2252F).  Problem NAMES are cited in prose / by \probref only; NO CF URL is
%  re-linked here (urlsweep: every URL lives in exactly one place -- 2127E/2252F
%  in the solvedproblems, 1292D in the Summit, the mixed ten in 99_chapter_end).
%  The opener's pattern card (00_opener, recognitionbox) is the one-glance
%  version; this section is the drill that reads five specs cold and routes them.
% =====================================================================
\section{Reading a Statement: Is It a Virtual Tree?}\label{sec:c30-recognition}

\begin{previouslybox}[The pattern card is up front; this is the repetition drill]
The opener handed you the one-line test -- \emph{one fixed tree, a marked set of size
\(k\ll n\), a computation repeated over that set} (\(\figref{c30-taxonomy}\)).  A
signature you have only read once is a signature you will misapply under contest
pressure.  This section does the boring, load-bearing thing: it takes five short
specifications -- three that are virtual trees and two confident traps -- and runs
each through the gate before a single line of code, so the reflex \emph{check the
marked set, then the tree} is automatic by the time you reach the mixed set.
Nothing here is solved to completion; the whole skill is the \emph{routing}.
\end{previouslybox}

\begin{topicintro}
A virtual tree is not a topic you recognise from a keyword -- statements never say
``compress the tree.'' They say ``for each colour, over the vertices of that
colour\ldots'' or ``given a set of \(k\) special nodes, compute\ldots'' and let an
\(O(n)\) bound sit next to a \(k\ll n\) bound so that re-walking the whole tree per
run is the thing that times out.  The recognition is arithmetic: count how many
times you would touch the full tree, and how many nodes actually matter per touch.
If the product \(n\times(\text{runs})\) is over budget but \(\sum(\text{marked})\)
is not, the tree has a \(\Theta(k)\) skeleton waiting.

\textbf{Prerequisites.} \(\secref{c30-construction}\) (what the skeleton is),
\(\secref{c30-dp}\) (what you run on it), \(\secref{c30-apps}\) (the two families).

\textbf{After this section you can}
\begin{itemize}[nosep]
  \item state the three-part virtual-tree signature as a yes/no gate; \hfill\lv{3}
  \item separate ``small set on a fixed tree'' from ``every node matters''; \hfill\lv{4}
  \item reject the traps -- whole-tree DP, graph Steiner, moving topology -- by name; \hfill\lv{4}
  \item price a run as \(\bigO{k\log k}\) versus \(\bigO{n}\) straight from the constraints. \hfill\lv{4}
\end{itemize}
\end{topicintro}

% ---------------------------------------------------------------
\subsection{The signature, as a three-question gate}
% ---------------------------------------------------------------

\begin{definition}{The virtual-tree recognition signature}{c30-signature}
A problem is a virtual-tree problem when \emph{all three} hold. \textbf{(S1) Fixed
tree.} One tree on \(n\) nodes is given once and is never edited (no link/cut) -- so
\(\tin,\tout,\depth,\lca\) are precomputed once and reused. \textbf{(S2) A small
marked set drives the answer.} The quantity asked is over the subgraph \emph{induced
by} a set \(S\) of \(k\) named vertices (a colour class, a query's terminals, a
subset) together with the paths connecting them -- the marked-to-marked Steiner
skeleton -- not over all \(n\) vertices. \textbf{(S3) The set is re-run, and the
total marked count is small.} Either there are many such sets whose sizes satisfy
\(\sum k \ll n\cdot(\text{runs})\), or a single set with \(k\ll n\) must be handled
faster than an \(O(n)\) walk. Then collapsing to \(A(S)\) with \(|A(S)|\le 2k-1\)
(\(\thmref{c30-size}\)) turns an \(O(n)\) per-run cost into \(O(k\log k)\).
If any question is \emph{no}, it is a different tool; the gate has a reason for
every rejection.
\end{definition}

The three questions are ordered on purpose.  S1 (fixed tree) is a hard veto: if the
topology moves you are in Chapter~25's territory, not here.  S2 (a marked-induced
answer) is the real discriminator -- it separates ``the DP is \emph{about} the set''
from ``the DP happens to run on a tree.''  S3 (the budget arithmetic) is what makes
the compression \emph{pay}, and is the one students skip and then rebuild the whole
tree for no reason.

\begin{warningbox}[The two ways the signature is faked]
\begin{itemize}[nosep,leftmargin=1.4em]
  \item \textbf{``Small \(k\)'' but the answer needs \emph{every node's} value.} A
        count of nodes \emph{not} in \(S\), or a value that depends on unmarked
        subtrees, cannot be read off \(A(S)\) -- skipped nodes carry no information
        you kept.  That is a full tree DP / rerooting (\chref{c10}, \chref{c20}); a
        virtual tree would drop the very nodes the answer depends on.
  \item \textbf{``Marked terminals on a tree'' but they may meet at a \emph{new}
        branch node.} On the \emph{given} tree the Steiner nodes are exactly the
        marked set plus their LCAs (\(\secref{c30-apps}\)); on a \emph{graph} the
        cheapest connector can be a non-terminal, so the union-of-paths picture
        breaks -- that is bitmask Steiner DP (\(\secref{c30-boundary}\)), not \(A(S)\).
\end{itemize}
\end{warningbox}

% ---------------------------------------------------------------
\subsection{Five specs, read and routed}
% ---------------------------------------------------------------

Each card is a genuine contest shape (the three positives are this chapter's own
solveds and Summit; the two traps are mixed-set entries).  Read the constraint line,
then answer S1/S2/S3 aloud, then name the tool.

\begin{itemize}[leftmargin=1.4em]
  \item \textbf{``\(n\le2\cdot10^5\) fixed rooted tree; each vertex has a colour in
        \([0,k]\); for each colour, reason over the vertices carrying it; repaint
        blanks to minimise a weight sum.''} S1 yes (one tree, no edits). S2 yes --
        the cutie test is about the \(\lca\)-branch points of each \emph{colour
        class}, exactly the internal nodes of that colour's virtual tree. S3 yes --
        \(\sum_c|S_c|=n\) but each colour is touched once, so a per-colour
        \(\bigO{|S_c|\log}\) beats \(\bigO{n}\) per colour. \(\Rightarrow\) \textbf{VT},
        one per colour (\(\probref{c30-ancient}\)).
  \item \textbf{``\(n\le2\cdot10^5\) tree; vertices coloured; for each colour pick a
        connected \(k_c\)-set minimising the sum of distances from that colour's
        vertices to it; \(\sum n\le2\cdot10^5\).''} S1 yes. S2 yes -- the answer
        routes every colour-\(c\) distance through the colour-\(c\) Steiner skeleton,
        so only that colour's virtual tree has non-zero edge weight. S3 yes -- a
        whole-tree DP per colour would be \(\bigO{n\cdot\#\text{colours}}\); the
        per-colour virtual trees cost \(\bigO{\sum_c |S_c|\log}=\bigO{n\log}\).
        \(\Rightarrow\) \textbf{VT + fold + greedy} (\(\probref{c30-spectral}\)).
  \item \textbf{``An infinite divisor tree; \(n\) fragments sit at factorial nodes
        \(k_i!\); choose an assembly node \(P\) minimising \(\sum\) path lengths;
        \(K!\le10^{18}\) so the reachable nodes are \(\le O(\log)\) deep chains.''}
        S1 yes (the divisor tree never changes). S2 yes -- only the factorial
        nodes and their LCAs (which are factorials) lie on any fragment-to-fragment
        path. S3 yes -- the factorials \(1!,2!,\dots,K!\) are already in \(\tin\)-
        order, so the marked set is closed and the virtual tree is \(\le 2n\) nodes
        while the real tree is astronomically large. \(\Rightarrow\) \textbf{VT +
        weighted centroid} -- the Summit (\(\probref{c30-summit}\)).
  \item \textbf{``A tree; every vertex has a value; answer questions that depend on
        the whole tree's aggregate, e.g.\ a sum over all \(n\) nodes / a diameter.''}
        S1 may hold, but S2 fails: the answer is \emph{about every node}, not the
        induced subgraph of a marked set -- there is no small set to collapse.
        \(\Rightarrow\) \emph{not} VT; tree DP / rerooting (\chref{c10}, \chref{c20}).
        This is the trap behind mixed entries like CF 1901E and CF 2253E: a big \(n\),
        a global answer, \emph{no} marked set.
  \item \textbf{``A weighted \emph{graph}; pick a minimum subgraph connecting \(k\)
        terminals; \(k\le 18\).''} S1 fails on the word \emph{graph} (the terminals
        connect through arbitrary intermediate vertices, possibly non-terminals), and
        the ``new branch node'' fake-out fires. \(\Rightarrow\) \emph{not} VT; bitmask
        Steiner DP (Dreyfus--Wagner), \(\bigO{3^k m + 2^k n\log n}\) -- the honest
        contrast in the mixed set is CF 152E. On a \emph{tree} this same question
        \emph{would} be free via \(A(S)\) (\(\secref{c30-apps}\)); the graph is the veto.
\end{itemize}

\begin{comparisonbox}[The same three questions, run on the two traps]
\setlength{\tabcolsep}{4pt}\renewcommand{\arraystretch}{1.2}%
\begin{tabular}{@{}p{0.30\linewidth}ccc@{}}
\toprule
\textbf{Spec} & \textbf{S1 fixed?} & \textbf{S2 set-induced?} & \textbf{S3 \(k\ll n\)?}\\
\midrule
per-colour cutie weights (\(\probref{c30-ancient}\)) & yes & yes & yes \(\Rightarrow\) VT \\
per-colour distance-to-component (\(\probref{c30-spectral}\)) & yes & yes & yes \(\Rightarrow\) VT \\
fragment assembly on the divisor tree (Summit) & yes & yes & yes \(\Rightarrow\) VT \\
global sum / diameter over all nodes & yes & \textbf{no} & -- \(\Rightarrow\) tree DP \\
min connector of \(k\) terminals in a \textbf{graph} & \textbf{no} & new nodes & yes \(\Rightarrow\) bitmask Steiner \\
\bottomrule
\end{tabular}
\end{comparisonbox}

\begin{insightbox}[The budget line is the honest signal, not the phrase ``marked set'']
Statements will describe ``special vertices'' without ever hinting at compression, and
they will describe a marked set that is \emph{the whole tree}.  The arithmetic is what
talks: write the naive cost \(n\cdot(\text{runs})\) and the collapse cost
\(\sum(\text{marked})\cdot\log\).  When the first blows the time limit and the second
does not, and the answer provably depends only on marked-to-marked paths, a virtual
tree is not a clever extra -- it is the only thing that fits.  When the second is
\(O(n)\) anyway (every node marked) the compression buys nothing and you should keep
the plain tree DP you already wrote.
\end{insightbox}

\begin{hardeningbox}
\begin{itemize}[nosep,leftmargin=1.4em]
  \item \textbf{Check S1 before anything.} If a link or cut appears, stop -- Chapter~25
        (\chref{c23}), not here; a virtual tree is a static-tree tool.
  \item \textbf{Ask ``does an unmarked node ever change the answer?''} If yes, S2 is
        false and you need the full tree; if no, you are entitled to drop every node
        outside the marked Steiner skeleton (\(\thmref{c30-fold}\)).
  \item \textbf{Tree vs graph Steiner is the sharpest fork.} Same words (``connect the
        terminals''), different object: on a tree the connectors are the marked set +
        their LCAs (free, \(A(S)\)); on a graph they are new nodes (hard, bitmask DP).
\end{itemize}
\end{hardeningbox}

\begin{recallbox}[26.4 in one breath]
\emph{Three questions, in order: is the tree fixed (else Ch.~25); does the answer live
on the marked set's induced subtree (else tree DP / reroot); is \(\sum k\) small enough
that \(O(k\log k)\) beats \(O(n)\) per run (else no gain).  All yes \(\Rightarrow\) build
\(A(S)\) and run the fold (\(\defref{c30-signature}\)).  The two traps -- ``the whole
tree's aggregate'' and ``terminals on a graph meet at new nodes'' -- are the reasons the
mixed set hands work to tree DP and bitmask Steiner, not to a virtual tree.}
\selftest{
  \item State the three questions of the gate. \emph{S1 fixed tree; S2 answer is marked-set-induced; S3 \(\sum k\) makes \(O(k\log k)\) beat \(O(n)\) (\(\defref{c30-signature}\)).}
  \item Why does a per-colour problem that recomputes an \(O(n)\) DP per colour want a virtual tree? \emph{Because \(\sum_c|S_c|=n\) and each colour's answer routes through its own \(A(S_c)\), so the total drops to \(O(n\log n)\).}
  \item Give the one-word veto that sends a ``connect the terminals'' problem away from a virtual tree. \emph{``graph'' -- terminals may meet at new non-terminal nodes; bitmask Steiner DP (\(\secref{c30-boundary}\)).}
}
\end{recallbox}
